% formal/search-pruning-correctness.tex — input by formal/main.tex
%
% Status: A-axis complete. Four variants: A0 (unpruned), A1 (vertex adjacency),
%   A2 (directed ω₀), A3 (Reeb-flow feasibility).
% Historical empirical ablation output was retired from the active tree; use git
% history for the old datasets, iteration counts, and timing artifacts if needed.
%
% Mathematical references used here (all Jörn-approved in commit 0839595):
%   lem:adjacency-constraint — consecutive facets with positive dwell time are adjacent
%   cor:adjacency-pruning — discard (S,σ) with non-adjacent consecutive pairs

\subsection{Pruning the Search Graph}
\label{sec:ablation-pruning}

We iteratively refine the adjacency graph used to prune the search space
of the HK2017 algorithm, starting from the unpruned baseline and adding
necessary conditions at each step; some refinements coincide on particular
polytope classes.

\paragraph{Variants.}
Table~\ref{tab:ablation-variants} lists the four adjacency-graph variants.
Each adds a necessary condition for the existence of a valid orbit transition
$F_i \to F_j$, building on all previous conditions.

\begin{table}[h]
\centering
\begin{tabular}{lp{9cm}}
\toprule
Variant & Pruning condition for transition $F_i \to F_j$ \\
\midrule
A0 & None (unpruned): exhaustive search over all $(S, \sigma)$ \\
A1 & Vertex adjacency: $F_i \cap F_j \neq \varnothing$, checked via shared polytope vertices (Corollary~\ref{cor:adjacency-pruning}) \\
A2 & Directed $\omega_0$: A1 and $\omega_0(a_i, a_j) \ge 0$, ensuring the Reeb flow on $F_i$ carries the orbit toward~$F_j$ \\
A3 & Reeb-flow feasibility: A2 and $\exists\, x \in F_i \cap F_j$, $\exists\,\varepsilon>0$ with $x - \varepsilon V_i \in F_i$ and $x + \varepsilon V_j \in F_j$, ensuring other halfspaces do not block the transition \\
\bottomrule
\end{tabular}
\caption{Adjacency graph variants, each at least as restrictive as the previous.
  The pruning condition applies to each consecutive pair in a candidate ordering
  $\sigma$; orderings that violate any edge condition are discarded.}
\label{tab:ablation-variants}
\end{table}

\paragraph{Directed symplectic condition (A2).}
For a physical orbit transition $F_i \to F_j$ at a ridge point $x \in F_i \cap F_j$,
the backward flow $x - \varepsilon V_i$ must remain in~$K$.
Checking the half-space constraint for facet~$j$ gives
$\omega_0(a_i, a_j) \ge 0$ as a necessary condition;
the forward flow $x + \varepsilon V_j$ on~$F_j$ gives the same condition
by antisymmetry of~$\omega_0$.
This is condition~(1) of Lemma~\ref{lem:transition-feasibility},
which proves the condition is both necessary and sufficient
(together with an LP feasibility check on the blocking set).

\paragraph{Reeb-flow feasibility (A3).}
A3 strengthens A2 by also checking the half-space constraints of all \emph{other}
facets $F_k$ ($k \neq i, j$).
For each $k$, we define a ``blocking set''
$B = \{k \notin \{i,j\} : \omega_0(a_i, a_k) < 0 \text{ or } \omega_0(a_j, a_k) > 0\}$,
and check whether $F_i \cap F_j$ contains a point satisfying
$\langle a_k, x \rangle < 1$ for all $k \in B$.
This is a strict linear feasibility problem on the face $F_i \cap F_j$.
A3 is a mathematical refinement in this account.  The active
\texttt{facet\_adjacency} implementation uses A2 (nonempty intersection
and the directed symplectic sign); it does not perform this blocker-slack
check.  The retired empirical ablation does not establish a current
implementation or runtime bound for A3.

\paragraph{Transition feasibility.}
We now characterize exactly when a piecewise-linear orbit can
transition from~$F_i$ to~$F_j$.
Write $V_k = J_0 a_k$ for the characteristic direction on~$F_k$
(proportional to the Reeb vector $R_k = 2\, J_0 a_k$
of Section~\ref{sec:algorithms}, but without the factor of~$2$).
In coordinates $(q_1, q_2, p_1, p_2)$ with
$J_0 = \bigl[\begin{smallmatrix} 0 & -I_2 \\ I_2 & 0 \end{smallmatrix}\bigr]$,
we have $V_k = (-a_k^{(3)}, -a_k^{(4)}, a_k^{(1)}, a_k^{(2)})$
and $a_k \cdot V_\ell = \omega_0(a_\ell, a_k)$.
A transition $F_i \to F_j$ requires a point $x \in F_i \cap F_j$
and $\varepsilon > 0$ with
$x + t\, V_i \in K$ for $t \in [-\varepsilon, 0]$
and $x + t\, V_j \in K$ for $t \in [0, \varepsilon]$.
The canonical lemma shows that the A3 feasibility check
captures this condition exactly.

The canonical statement and proof are
Lemma~\ref{lem:numerical-transition-feasibility}; the label
\texttt{lem:transition-feasibility} is an alias on that same lemma.
The hypotheses and blocking set above are identical.  In particular, when
$B=\emptyset$, the strict inequalities are vacuous; the face
$F_i\cap F_j$ must still be nonempty.  The proof chooses any positive
travel time when every directional constraint derivative vanishes.


\begin{unverified}
\begin{corollary}[Ridge sufficiency]
\label{cor:ridge-sufficiency}
If\/ $F_i \cap F_j$ is a 2-dimensional face (ridge) of~$K$
and $\omega_0(a_i, a_j) \ge 0$, then the transition $F_i \to F_j$
is feasible.
In particular, the A2 and~A3 conditions are equivalent
for simple polytopes and for Lagrangian products of convex polygons.
\end{corollary}

\begin{proof}
If $F_i \cap F_j$ is a ridge, any $x$ in its relative interior satisfies
$\langle a_k, x \rangle < 1$ for all $k \notin \{i,j\}$ (hence for all $k \in B$),
so condition~(2) of Lemma~\ref{lem:transition-feasibility} holds
automatically.
For simple 4-polytopes, each vertex lies on exactly four facets,
so any two adjacent facets share a ridge.
For Lagrangian products $P \times Q$ of convex polygons,
each vertex $(v_P, v_Q)$ lies on exactly two edges of~$P$
and two edges of~$Q$, hence on exactly four facets;
thus $P \times Q$ is simple, and every pair of intersecting facets
shares a 2-face (``vertex $\times$ polygon'' or ``edge $\times$ edge'').
\end{proof}
\end{unverified}

\begin{example}[Non-simple polytope where A3 prunes beyond A2]
\label{ex:a3-prunes}
Let $S$ be the 4-simplex with vertices
$v_0 = (2,0,0,0)$,\; $v_1 = (0,2,0,0)$,\; $v_2 = (0,0,2,0)$,\;
$v_3 = (0,0,0,2)$,\; $v_4 = (-2,-2,-2,-2)$
(centroid at the origin), and set
$K = S \cap \{\, x_1 + 2\, x_2 \le 2 \,\}$.
The cutting hyperplane passes through~$v_0$ and removes~$v_1$,
producing a polytope with $6$~facets and $7$~vertices.
The vertex~$v_0$ lies on $5$~facets ($F_1,\ldots,F_4$ from the simplex
plus the cutting facet~$F_5$), so $K$ is non-simple.
Every facet pair shares a ridge \emph{except} $(F_1, F_5)$,
whose intersection is the single vertex~$v_0$.

Since $\omega_0(a_5, a_1) > 0$, condition~(1) of
Lemma~\ref{lem:transition-feasibility} holds and A2 allows
the directed transition $F_5 \to F_1$.
However, $F_2$ belongs to the blocking set~$B$
(since $\omega_0(a_5, a_2) < 0$), and $v_0 \in F_2$
gives $\langle a_2, v_0 \rangle = 1$:
the only point of $F_1 \cap F_5$ is tight on the blocking facet.
Condition~(2) fails, and A3 correctly prunes this transition.
The retired ablation reported $33$ candidate orderings for A3 versus $39$
for~A2; these counts were not regenerated in the source consolidation.
\end{example}
